\documentclass[10pt]{article}

\usepackage[margin=0.8in]{geometry}
\usepackage{amsmath,amssymb,booktabs}
\usepackage{algorithm}
\usepackage{algpseudocode}
\usepackage{microtype}
\usepackage{enumitem}
\usepackage[hidelinks]{hyperref}

\title{GAP-DLM: Implementation Notes for Retain-Free Unlearning\\
in Masked Diffusion Language Models}
\author{}
\date{}

\begin{document}
\maketitle

\section{Goal}

We assume a pretrained masked diffusion language model (MDLM) and a forget
set only. The goal is to remove prompt-dependent knowledge associated with
forget examples while preserving the original model behavior as much as possible.

The method uses a frozen copy of the original model as a teacher. For each
masked denoising state, we compare the teacher prediction conditioned on the
forget prompt against the teacher prediction with the prompt removed. We only
suppress probability mass that is specifically amplified by the forget prompt.
All other behavior is self-distilled from the original model.

\section{Notation}

\begin{table}[h]
\centering
\small
\begin{tabular}{@{}ll@{}}
\toprule
Symbol & Meaning \\
\midrule
$\theta_0$ & Frozen original MDLM teacher \\
$\theta$ & Trainable unlearned MDLM \\
$\mathcal{D}_f$ & Forget dataset of pairs $(x,y)$ \\
$x$ & Forget prompt / input \\
$y=(y_1,\dots,y_L)$ & Target response \\
$t$ & Diffusion noise / mask level \\
$y_t$ & Partially masked response at level $t$ \\
$\mathcal{M}_t$ & Masked token positions in $y_t$ \\
$m$ & Null prompt \\
$p^c_{t,i}$ & Teacher distribution conditioned on $(x,y_t)$ \\
$p^u_{t,i}$ & Teacher distribution conditioned on $(m,y_t)$ \\
$p^\theta_{t,i}$ & Student distribution conditioned on $(x,y_t)$ \\
$g_{t,i}(v)$ & Prompt-induced log-probability lift for token $v$ \\
$\kappa_i$ & Copy-safe indicator for target token $y_i$ \\
$\alpha_{t,i}$ & Forget strength at token-state pair $(t,i)$ \\
$\mathcal{S}^{+}_{t,i}$ & Tokens whose probability is increased by the prompt \\
$q^*_{t,i}$ & Projected teacher target \\
$\lambda$ & Maximum unlearning strength in $[0,1]$ \\
$\epsilon$ & Small constant for stable logarithms \\
\bottomrule
\end{tabular}
\caption{Main notation.}
\end{table}

\section{Teacher Distributions}

For each forget example $(x,y)$, sample a mask level $t$ and construct a
partially masked response $y_t$. For every masked position $i\in\mathcal{M}_t$,
compute
\begin{align}
p^c_{t,i}(v) &= p_{\theta_0}(v\mid x,y_t),\\
p^u_{t,i}(v) &= p_{\theta_0}(v\mid m,y_t).
\end{align}

The conditional distribution $p^c$ is the original model behavior on the
forget example. The null distribution $p^u$ estimates what can be predicted
from the partially observed response without using the forget prompt.

\section{Prompt-Attributed Guidance}

Define the prompt-induced log-probability lift
\begin{equation}
g_{t,i}(v)=
\log\left(p^c_{t,i}(v)+\epsilon\right)
-
\log\left(p^u_{t,i}(v)+\epsilon\right).
\end{equation}

A positive value means that the forget prompt increases the probability of
token $v$. We use the ground-truth token $y_i$ to decide how strongly the
current token-state pair should be modified:
\begin{equation}
\alpha_{t,i}
=
\lambda\kappa_i
\left(1-\exp\left(-[g_{t,i}(y_i)]_+\right)\right),
\end{equation}
where $[z]_+=\max(z,0)$.

When $p^c_{t,i}(y_i)>p^u_{t,i}(y_i)$, this is equivalently
\begin{equation}
\alpha_{t,i}
=
\lambda\kappa_i
\left(
1-\frac{p^u_{t,i}(y_i)}{p^c_{t,i}(y_i)}
\right).
\end{equation}

If the prompt does not help predict the target token, then
$\alpha_{t,i}\approx0$ and the original behavior is preserved.

\subsection{Copy-Safe Indicator}

We do not want to erase the ability to copy information explicitly provided in
the prompt. Define
\begin{equation}
\kappa_i=
\begin{cases}
0, & \text{if $y_i$ belongs to a span copied from $x$,}\\
1, & \text{otherwise.}
\end{cases}
\end{equation}

A simple implementation can use exact token-span matching or longest common
subsequence matching between $x$ and $y$.

\section{Positive-Guidance Belief Set}

Instead of suppressing only the labeled target token, suppress all vocabulary
tokens whose probability is increased by the forget prompt:
\begin{equation}
\mathcal{S}^{+}_{t,i}
=
\{v:g_{t,i}(v)>0\}.
\end{equation}

This reduces the chance that probability mass moves from the original target
token to an alternative token carrying similar information.

\section{Projected Teacher Target}

For every token $v\in\mathcal{S}^{+}_{t,i}$, define an upper bound
\begin{align}
u_{t,i}(v)
&=
p^c_{t,i}(v)\exp\left(-\alpha_{t,i}g_{t,i}(v)\right)\\
&=
\left(p^c_{t,i}(v)\right)^{1-\alpha_{t,i}}
\left(p^u_{t,i}(v)\right)^{\alpha_{t,i}}.
\end{align}

The target distribution is the closest distribution to the original
conditional teacher while respecting these upper bounds:
\begin{equation}
q^*_{t,i}
=
\arg\min_{q\in\Delta^{|\mathcal{V}|}}
\mathrm{KL}\left(q\,\|\,p^c_{t,i}\right)
\end{equation}
subject to
\begin{equation}
q(v)\le u_{t,i}(v),
\qquad
\forall v\in\mathcal{S}^{+}_{t,i}.
\end{equation}

This has the closed-form solution
\begin{equation}
q^*_{t,i}(v)=
\begin{cases}
u_{t,i}(v), & v\in\mathcal{S}^{+}_{t,i},\\[4pt]
\gamma_{t,i}p^c_{t,i}(v), & v\notin\mathcal{S}^{+}_{t,i},
\end{cases}
\end{equation}
where
\begin{equation}
\gamma_{t,i}
=
\frac{
1-\sum_{v\in\mathcal{S}^{+}_{t,i}}u_{t,i}(v)
}{
\sum_{v\notin\mathcal{S}^{+}_{t,i}}p^c_{t,i}(v)
}.
\end{equation}

Important special case:
\begin{equation}
\alpha_{t,i}=0
\quad\Longrightarrow\quad
q^*_{t,i}=p^c_{t,i}.
\end{equation}
Thus, token-state pairs without clear forget evidence become ordinary
self-distillation examples.

\section{Training Objective}

The student is trained with one KL-divergence objective:
\begin{equation}
\mathcal{L}_{\mathrm{GAP}}
=
\mathbb{E}_{(x,y),t,y_t}
\left[
\frac{1}{|\mathcal{M}_t|}
\sum_{i\in\mathcal{M}_t}
\mathrm{KL}
\left(
\operatorname{sg}[q^*_{t,i}]
\;\middle\|\;
p^\theta_{t,i}
\right)
\right],
\end{equation}
where
\begin{equation}
p^\theta_{t,i}=p_\theta(\cdot\mid x,y_t).
\end{equation}
The projected teacher target is detached before the student update.

\section{Training Algorithm}

\begin{algorithm}[h]
\caption{GAP-DLM Retain-Free Unlearning}
\begin{algorithmic}[1]
\Require Original model $\theta_0$, forget set $\mathcal{D}_f$, strength $\lambda$
\State Initialize student $\theta\leftarrow\theta_0$
\State Freeze $\theta_0$
\For{each training step}
    \State Sample minibatch $B\subset\mathcal{D}_f$
    \For{each $(x,y)\in B$}
        \State Sample mask level $t$ and construct $y_t$
        \State Compute copy indicators $\kappa_i$
        \State Compute teacher distributions
        \[
        p^c=p_{\theta_0}(\cdot\mid x,y_t),
        \qquad
        p^u=p_{\theta_0}(\cdot\mid m,y_t)
        \]
        \For{each masked position $i\in\mathcal{M}_t$}
            \State $g_i\leftarrow\log(p^c_i+\epsilon)-\log(p^u_i+\epsilon)$
            \State $\alpha_i\leftarrow
            \lambda\kappa_i(1-\exp(-[g_i(y_i)]_+))$
            \State $\mathcal{S}^+_i\leftarrow\{v:g_i(v)>0\}$
            \State $u_i(v)\leftarrow
            (p^c_i(v))^{1-\alpha_i}(p^u_i(v))^{\alpha_i}$
            for $v\in\mathcal{S}^+_i$
            \State Construct projected target $q^*_i$
        \EndFor
        \State Compute student distribution $p^\theta=p_\theta(\cdot\mid x,y_t)$
        \State Accumulate
        $\frac{1}{|\mathcal{M}_t|}\sum_{i\in\mathcal{M}_t}
        \mathrm{KL}(q^*_i\|p^\theta_i)$
    \EndFor
    \State Update $\theta$ with gradient descent
\EndFor
\State \Return unlearned model $\theta$
\end{algorithmic}
\end{algorithm}

\section{Minimal Vectorized Implementation}

For numerical stability, operate in log-space. Given conditional and null
teacher logits $z^c$ and $z^u$:
\begin{align}
\ell^c &= \log\operatorname{softmax}(z^c),\\
\ell^u &= \log\operatorname{softmax}(z^u),\\
g &= \ell^c-\ell^u.
\end{align}

Then compute
\begin{align}
\texttt{positive} &= (g>0),\\
\alpha &= \lambda\kappa
\left(1-\exp(-\max(g_{\mathrm{target}},0))\right),\\
\log u &= (1-\alpha)\ell^c+\alpha\ell^u.
\end{align}

Only masked response positions should contribute to the loss. A direct
full-vocabulary implementation is recommended first. If memory becomes a
bottleneck, the projection can later be approximated with the union of the
top-$K$ tokens from $p^c$ and $p^u$.

\section{Recommended First Implementation}

Use the simplest version first:
\begin{itemize}[leftmargin=1.5em]
    \item backbone: LLaDA or Dream;
    \item one frozen teacher and one trainable student;
    \item the same masking distribution used by the MDLM training code;
    \item exact token-span matching for $\kappa_i$;
    \item $\epsilon=10^{-8}$;
    \item $\lambda\in\{0.25,0.5,0.75,1.0\}$;
    \item full-vocabulary projection;
    \item no hard-mask mining and no extra regularizer.
\end{itemize}

This version isolates whether conditional--null guidance is enough to identify
forget-relevant denoising decisions before adding more complex components.

\section{Sanity Checks}

Before full training, verify:
\begin{enumerate}[leftmargin=1.5em]
    \item If $p^c(y_i)\approx p^u(y_i)$, then $\alpha_{t,i}\approx0$.
    \item If $\kappa_i=0$, then $q^*_{t,i}=p^c_{t,i}$.
    \item If $\alpha_{t,i}=0$, the loss is ordinary self-distillation.
    \item For every $v\in\mathcal{S}^+_{t,i}$, $q^*_{t,i}(v)\le u_{t,i}(v)$.
    \item Every projected target satisfies $\sum_v q^*_{t,i}(v)=1$.
    \item The target $q^*$ is detached before backpropagation.
\end{enumerate}

\end{document}
